\section{局部化与 $n$-同调}
\label{sec:03.3.2}

设 $\mathcal{M}(U(\g))$ 是 $U(\g)$-模构成的范畴。固定一点 $x\in X$。对 $V\in\mathcal{M}(U(\g))$，置
$$V_{\n_x}=V/\n_xV=\C\otimes_{U(\n_x)}V,$$
其中把 $\C$ 看作带 $\n_x$ 平凡作用的模。称 $V_{\n_x}$ 为 $V$ 中的 $\n_x$-余不变量模。它有自然的 $\h_x$-模结构，因此可看作 $\h$-模。由此可见，$V\mapsto V_{\n_x}$ 是从范畴 $\mathcal{M}(U(\g))$ 到 $U(\h)$-模范畴 $\mathcal{M}(U(\h))$ 的右正合共变函子。若再与从 $\mathcal{M}(U(\h))$ 到向量空间范畴的遗忘函子复合，就得到零次 $\n_x$-同调函子 $H_0(\n_x,-)$。由 Poincaré--Birkhoff--Witt 定理，自由 $U(\g)$-模也是自由 $U(\n_x)$-模，故二者的左导出函子相同。因此，术语上稍有滥用，我们把 $V\mapsto V_{\n_x}$ 的第 $(-p)$ 个左导出函子称为第 $p$ 个 $\n_x$-同调函子，并记作
$$H_p(\n_x,-)=\Tor_p^{U(\n_x)}(\C,-).$$
这些函子与局部化函子之间有简单的关系，下面说明。

先需要一个技术性结果。

\begin{lemma}\label{lem:03.2.1}
$U_\theta$ 作为 $U(\n_x)$-模是自由的。
\end{lemma}

\begin{proof}
固定 $\h$ 的一个特化 $\mathfrak{c}$，并取与 $\n_x$ 相对的幂零子代数 $\bar{\n}$。则 $\g=\n_x\oplus\mathfrak{c}\oplus\bar{\n}$；由 Poincaré--Birkhoff--Witt 定理可知，作为左 $U(\n_x)$-模（对左乘），
$$U(\g)=U(\n_x)\otimes_\C U(\mathfrak{c})\otimes_\C U(\bar{\n}).$$
设 $F_pU(\mathfrak{c})$（$p\in\Z_+$）是 $U(\mathfrak{c})$ 的次数滤过。于是通过
$$F_pU(\g)=U(\n_x)\otimes_\C F_pU(\mathfrak{c})\otimes_\C U(\bar{\n})$$
定义 $U(\g)$ 的滤过 $F_pU(\g)$（$p\in\Z_+$）。这显然是 $U(\n_x)$-模滤过。相应的分次模为
$$\Gr U(\g)=U(\n_x)\otimes_\C S(\mathfrak{c})\otimes_\C U(\bar{\n}).$$
这个滤过在子模 $J_\theta U(\g)$ 与商模 $U_\theta$ 上诱导滤过。Harish-Chandra 同态 $\gamma:Z(\g)\to U(\h)$ 与次数滤过相容，且同态 $\Gr\gamma$ 是 $\Gr Z(\g)$ 到 $S(\h)$ 中全体 $W$-不变元所成子代数 $I(\h)$ 的同构 [5, Ch.~VIII, \S8, no.~5]。记 $I_+(\h)$ 为由 $I(\h)$ 中严格正次数的元素生成的齐次理想。则
$$\Gr J_\theta U(\g)=U(\n_x)\otimes_\C I_+(\mathfrak{c})S(\mathfrak{c})\otimes_\C U(\bar{\n}).$$
由 [4, Ch.~III, \S2, no.~4, Prop.~2] 可得
\begin{align*}
\Gr U_\theta&=\frac{\Gr U(\g)}{\Gr J_\theta U(\g)}\\
&=\frac{U(\n_x)\otimes_\C S(\mathfrak{c})\otimes_\C U(\bar{\n})}
{U(\n_x)\otimes_\C I_+(\mathfrak{c})S(\mathfrak{c})\otimes_\C U(\bar{\n})}\\
&=U(\n_x)\otimes_\C\frac{S(\mathfrak{c})}{I_+(\mathfrak{c})S(\mathfrak{c})}\otimes_\C U(\bar{\n}),
\end{align*}
即它是自由 $U(\n_x)$-模。此外，由 [5, Ch.~V, \S5, no.~2, Th.~1] 知，复向量空间 $S(\h)/(I_+(\h)S(\h))$ 的维数为 $\operatorname{Card}W$。由此可见，$U_\theta$ 有一个由 $U(\n_x)$-子模构成的有限滤过，使 $\Gr U_\theta$ 是自由 $U(\n_x)$-模。对滤过长度归纳，即知 $U_\theta$ 是自由 $U(\n_x)$-模。
\end{proof}

设 $\varphi:U(\h)\to U(\h)$ 是由 $\varphi(\xi)=\xi+\rho(\xi)$（$\xi\in\h$）给出的自同构。则由 [5, Ch.~VIII, \S8, no.~5, Th.~2] 知，$\varphi(\gamma(Z(\g)))$ 是 $U(\h)$ 中 $W$-不变元构成的代数。此外，由 [5, Ch.~V, \S5, no.~2, Th.~1] 知，向量空间 $U(\h)/\varphi(\gamma(J_\theta))U(\h)$ 的维数等于 $\operatorname{Card}W$。由此推出，$V_\theta=U(\h)/\gamma(J_\theta)U(\h)$ 是维数为 $\operatorname{Card}W$ 的 $U(\h)$-模。

\begin{lemma}\label{lem:03.2.2}
设 $\lambda\in\h^*$，$\theta=W\cdot\lambda$。则：
\begin{enumerate}[label=(\roman*)]
\item $V_\theta$ 是维数为 $\operatorname{Card}W$ 的 $U(\h)$-模；
\item $\xi\in\h$ 在 $V_\theta$ 上作用的特征多项式为
$$P(\xi)=\prod_{w\in W}\bigl(\xi-(w\lambda+\rho)(\xi)\bigr);$$
\item $H_0(\n_x,U_\theta)$ 是可数多个 $V_\theta$ 的直和。
\end{enumerate}
\end{lemma}

\begin{proof}
(i) 已证。(ii) 显然，$I_\mu\supset\varphi(\gamma(J_\theta))U(\h)$ 等价于存在 $w\in W$ 使 $\mu=w\lambda$。因此，由乘以 $\xi$ 在 $U(\h)/\varphi(\gamma(J_\theta))U(\h)$ 上诱导的线性变换有特征值 $(w\lambda)(\xi)$（$w\in W$），且由对称性它们的重数\index{chongshu@重数|hyperpage}都相同。这又推出
$$\varphi(P(\xi))=\prod_{w\in W}\varphi\bigl(\xi-(w\lambda+\rho)(\xi)\bigr)=\prod_{w\in W}\bigl(\xi-(w\lambda)(\xi)\bigr)$$
是 $\xi$ 在 $U(\h)/\varphi(\gamma(J_\theta))U(\h)$ 上作用的特征多项式。这就证明了 (ii)。

(iii) 与 \ref{lem:03.2.1} 的证明一样，固定 $\h$ 的一个特化 $\mathfrak{c}$，并取与 $\n_x$ 相对的幂零子代数 $\bar{\n}$。由 Poincaré--Birkhoff--Witt 定理可知，作为向量空间 $U(\g)=U(\n_x)\otimes_\C U(\mathfrak{c})\otimes_\C U(\bar{\n})$。此外
$$H_0(\n_x,U_\theta)=U(\g)/(J_\theta U(\g)+\n_xU(\g)).$$
记 $\gamma_x:Z(\g)\to U(\mathfrak{c})$ 为特化映射与 Harish-Chandra 同态 $\gamma$ 的复合。则
$$J_\theta U(\g)+\n_xU(\g)=J_\theta U(\mathfrak{c})U(\bar{\n})+\n_xU(\g)=\gamma_x(J_\theta)U(\mathfrak{c})U(\bar{\n})+\n_xU(\g),$$
由此可见，在上述同构下
\begin{align*}
J_\theta U(\g)+\n_xU(\g)
&=\bigl(U(\n_x)\otimes_\C\gamma_x(J_\theta)U(\mathfrak{c})\otimes_\C U(\bar{\n})\bigr)\\
&\qquad\oplus\bigl(\n_xU(\n_x)\otimes_\C U(\mathfrak{c})\otimes_\C U(\bar{\n})\bigr).
\end{align*}
\begin{annote}
底本此处把第一个因子印作 $\C$。但紧随其后的商计算把零次同调算成 $U(\mathfrak{c})/\gamma_x(J_\theta)U(\mathfrak{c})$ 与 $U(\bar{\n})$ 的张量积；而等式右端已有 $\n_xU(\n_x)\otimes_\C U(\mathfrak{c})\otimes_\C U(\bar{\n})$ 一项，故第一个因子中尚需保留的正是整个 $U(\n_x)$。已按上下文改正。[来源：底本 p.44，2.2(iii) 的证明，与紧接其后的商计算对照]
\end{annote}
由此得到
$$H_0(\n_x,U_\theta)=U(\mathfrak{c})/(\gamma_x(J_\theta)U(\mathfrak{c}))\otimes_\C U(\bar{\n})=V_\theta\otimes_\C U(\bar{\n}),$$
而 $\h$ 的作用由第一个因子中的乘法给出。
\end{proof}

\begin{corollary}\label{cor:03.2.3}
设 $\lambda\in\h^*$，$\theta=W\cdot\lambda$，$V\in\mathcal{M}(U_\theta)$。若置
$$P(\xi)=\prod_{w\in W}\bigl(\xi-(w\lambda+\rho)(\xi)\bigr),\qquad \xi\in\h,$$
则对任意 $\xi\in\h$ 与 $p\in\Z_+$，$P(\xi)$ 消灭 $H_p(\n_x,V)$。
\end{corollary}

\begin{proof}
由 \ref{lem:03.2.1}，可用 $V$ 的自由 $U_\theta$-模左消解来计算 $V$ 的 $\n_x$-同调。断言由 \ref{lem:03.2.2} 得出。
\end{proof}

特别地，若 $V\in\mathcal{M}(U_\theta)$，则 $H_p(\n_x,V)$ 作为 $U(\h)$-模是相应于 $w\lambda+\rho$（$w\in W$）的广义特征空间的直和。若 $U$ 是 $U(\h)$-模，对 $\lambda\in\h^*$，记 $U_{(\lambda)}$ 为相应于 $\lambda$ 的广义特征空间。

\begin{corollary}\label{cor:03.2.4}
设 $\lambda\in\h^*$ 正则。则对 $V\in\mathcal{M}(U_\theta)$，$\n_x$-同调模 $H_p(\n_x,V)$ 作为 $U(\h)$-模是半单的。更确切地说，对任意 $p\in\Z_+$ 有
$$H_p(\n_x,V)=\sum_{w\in W}H_p(\n_x,V)_{(w\lambda+\rho)}.$$
\end{corollary}

由此特别推出，对正则的 $\lambda\in\h^*$，可把函子 $H_p(\n_x,-)_{(w\lambda+\rho)}$ 看作从 $\mathcal{M}(U_\theta)$ 到 $\mathcal{M}(U(\h))$ 的右正合函子 $H_0(\n_x,-)_{(w\lambda+\rho)}$ 的第 $p$ 个左导出函子。

在一般情形，可把 $V_\theta$ 看作半局部环，并把 $H_p(\n_x,V)$ 看作 $V_\theta$-模。此外，对任意 $\lambda\in\theta$，$\C_{\lambda+\rho}$ 都是 $V_\theta$-模。

对 $X$ 上任意 $\calO_X$-模 $\mathcal{F}$，记 $T_x(\mathcal{F})$ 为其几何纤维，即
$$T_x(\mathcal{F})=\mathcal{F}_x/\m_x\mathcal{F}_x.$$
于是 $T_x$ 是从 $\mathcal{M}(\calO_X)$ 到复向量空间范畴的右正合共变函子。若 $\mathcal{F}$ 是 $\D_\lambda$-模，则可把 $T_x(\mathcal{F})$ 看作 $\mathcal{F}$ 沿单点空间 $\{x\}$ 到 $X$ 的包含映射 $i$ 的逆像\index{nixiang@逆像|hyperpage}。

\begin{theorem}\label{thm:03.2.5}
设 $\lambda\in\h^*$，$\theta=W\cdot\lambda$，$x\in X$。则从 $D^-(U_\theta)$ 到复向量空间复形的导出范畴的函子
$$LT_x\circ L\Delta_\lambda
\quad\text{与}\quad
D(\C_{\lambda+\rho})\overset{L}{\otimes}_{V_\theta}\bigl(D(\C)\overset{L}{\otimes}_{U(\n_x)}-\bigr)$$
同构。
\end{theorem}

\begin{proof}
由 \ref{lem:03.2.1} 知，$U_\theta$ 对函子 $H_0(\n_x,-)=\C\otimes_{U(\n_x)}-$ 零调。由 \ref{lem:03.2.2} 又知，$\C\otimes_{U(\n_x)}U_\theta$ 对函子 $\C_{\lambda+\rho}\otimes_{V_\theta}-$ 零调。设 $F^\bullet$ 是与 $V^\bullet$ 同构的、由自由 $U_\theta$-模构成的复形。于是，由于这些函子与无限直和交换，得到
$$D(\C_{\lambda+\rho})\overset{L}{\otimes}_{V_\theta}\bigl(D(\C)\overset{L}{\otimes}_{U(\n_x)}V^\bullet\bigr)=\C_{\lambda+\rho}\otimes_{V_\theta}\bigl(\C\otimes_{U(\n_x)}F^\bullet\bigr).$$
另一方面，局部化 $\Delta_\lambda(U_\theta)=\D_\lambda$ 是局部自由 $\calO_X$-模，因此对 $T_x$ 零调。这意味着
$$LT_x(L\Delta_\lambda(V^\bullet))=T_x(\Delta_\lambda(F^\bullet)).$$
于是，为完成证明，只需建立下述等式
$$T_x(\Delta_\lambda(U_\theta))=\C_{\lambda+\rho}\otimes_{V_\theta}\bigl(\C\otimes_{U(\n_x)}U_\theta\bigr).$$
首先，$T_x(\Delta_\lambda(U_\theta))=T_x(\D_\lambda)$。此外，由 $\D_\h$ 的构造可得
$$T_x(\D_\h)=U(\g)/\n_xU(\g),$$
再利用 Harish-Chandra 同态的性质，即得
\begin{align*}
T_x(\D_\lambda)&=T_x(\D_\h)\otimes_{U(\h)}\C_{\lambda+\rho}\\
&=\bigl(U(\g)/\n_xU(\g)\bigr)\big/\bigl(I_{\lambda+\rho}(U(\g)/\n_xU(\g))\bigr)\\
&=\C_{\lambda+\rho}\otimes_{V_\theta}\bigl(U(\g)/\n_xU(\g)\bigr)\big/\bigl(\gamma(J_\theta)(U(\g)/\n_xU(\g))\bigr)\\
&=\C_{\lambda+\rho}\otimes_{V_\theta}\bigl(U(\g)/(J_\theta U(\g)+\n_xU(\g))\bigr)
=\C_{\lambda+\rho}\otimes_{V_\theta}H_0(\n_x,U_\theta).
\end{align*}
\end{proof}

\begin{corollary}\label{cor:03.2.6}
设 $\lambda\in\h^*$ 正则，$\theta=W\cdot\lambda$。则对任意 $V\in\mathcal{M}(U_\theta)$ 有谱序列
$$L_pT_x\bigl(L_q\Delta_\lambda(V)\bigr)\Longrightarrow H_{-(p+q)}(\n_x,V)_{(\lambda+\rho)}.$$
\end{corollary}

\begin{proof}
如前所述，此时所有 $V_\theta$-模都半单，$\C_{\lambda+\rho}\otimes_{V_\theta}-$ 是正合函子。
\begin{annote}
底本此处印作 $\C_\lambda\otimes_{V_\theta}-$。但同一证明下文的计算用的是 $\C_{\lambda+\rho}$，且 $\C_\lambda$ 并不是 $V_\theta$-模（$V_\theta$ 的单模是 $\C_{w(\lambda+\rho)}$，$w\in W$）。故按上下文改为 $\C_{\lambda+\rho}\otimes_{V_\theta}-$。[来源：底本 p.46，2.6 的证明，与同段下文的计算对照]
\end{annote}
因此，与 \ref{thm:03.2.5} 中第二个函子相应的谱序列退化，得到
\begin{align*}
H^{-p}\bigl(D(\C_{\lambda+\rho})\overset{L}{\otimes}_{V_\theta}(D(\C)\overset{L}{\otimes}_{U(\n_x)}D(V))\bigr)
&=H^{-p}\bigl(\C_{\lambda+\rho}\otimes_{V_\theta}(D(\C)\overset{L}{\otimes}_{U(\n_x)}D(V))\bigr)\\
&=H^{-p}\bigl((D(\C)\overset{L}{\otimes}_{U(\n_x)}D(V))_{(\lambda+\rho)}\bigr)\\
&=H^{-p}\bigl(D(\C)\overset{L}{\otimes}_{U(\n_x)}D(V)\bigr)_{(\lambda+\rho)}\\
&=H_p(\n_x,V)_{(\lambda+\rho)}
\end{align*}
对 $p\in\Z_+$ 成立。因此，所断言的谱序列正是附着于 $LT_x$ 与 $L\Delta_\lambda$ 的复合的 Grothendieck 谱序列。
\end{proof}

奇异 $\lambda$ 处的行为要晦涩得多，由下述结果可见。

\begin{proposition}\label{prop:03.2.7}
设 $\lambda\in\h^*$ 奇异。则存在 $V\in\mathcal{M}(U_\theta)$，使 $L\Delta_\lambda(D(V))$ 不是有界复形。特别地，$\Delta_\lambda$ 的左上同调维数无穷。
\end{proposition}

\begin{proof}
由于函子 $T_x$ 的左上同调维数有限，只需找出一个 $U_\theta$-模 $V$，使对某个 $x\in X$，$LT_x(L\Delta_\lambda(V))$ 不是有界复形。由 \ref{thm:03.2.5}，这等价于
$$D(\C_{\lambda+\rho})\overset{L}{\otimes}_{V_\theta}\bigl(D(\C)\overset{L}{\otimes}_{U(\n_x)}D(V)\bigr)$$
不是有界复形。

为完成证明，我们用到关于 Verma 模的一些初等结果，它们将在后面的 V.1 中讨论。固定一个 Borel 子代数 $\bb_0$，置 $\n_0=[\bb_0,\bb_0]$，并考虑 Verma 模
$$M(w_0\lambda)=U(\g)\otimes_{U(\bb_0)}\C_{w_0\lambda-\rho}.$$
取 $x$，使 $\n_x$ 与 $\n_0$ 相对。则由 Poincaré--Birkhoff--Witt 定理，作为 $U(\n_x)$-模，$M(w_0\lambda)$ 同构于 $U(\n_x)\otimes_\C\C_{w_0\lambda-\rho}$。由此可知
$$H_0(\n_x,M(w_0\lambda))=\C_{\lambda+\rho},$$
且对 $p\in\N$，$H_p(\n_x,M(w_0\lambda))=0$。
\begin{annote}
底本由 $M(w_0\lambda)$ 作为 $U(\n_x)$-模自由，直接得出 $H_0(\n_x,M(w_0\lambda))=\C_{\lambda+\rho}$。但秩 $1$ 的自由模，其零次同调只是 $\C_{w_0\lambda-\rho}$；两者相等需 $w_0\lambda=\lambda+2\rho$，一般并不成立。故此处照录存疑，待考。[来源：底本 p.46，2.7 的证明；推理自证]
\end{annote}
因此
$$D(\C)\overset{L}{\otimes}_{U(\n_x)}D(M(w_0\lambda))=D(\C_{\lambda+\rho}),$$
并且
$$D(\C_{\lambda+\rho})\overset{L}{\otimes}_{V_\theta}\bigl(D(\C)\overset{L}{\otimes}_{U(\n_x)}D(M(w_0\lambda))\bigr)=D(\C_{\lambda+\rho})\overset{L}{\otimes}_{V_\theta}D(\C_{\lambda+\rho}).$$
显然
$$H^{-p}\bigl(D(\C_{\lambda+\rho})\overset{L}{\otimes}_{V_\theta}D(\C_{\lambda+\rho})\bigr)=\Tor_p^{V_\theta}(\C_{\lambda+\rho},\C_{\lambda+\rho}),\qquad p\in\Z_+.$$
设 $W(\lambda)$ 是 $\lambda$ 在 $W$ 中的稳定化子。由 \ref{lem:03.2.2}，$V_\theta$ 中的极大理想是理想 $I_{w\lambda+\rho}$ 的投影，$w\in W/W(\lambda)$。由于 $V_\theta$ 是 Artin 环，它是局部环 $R_{w\lambda}$（$w\in W/W(\lambda)$）的积，这些局部环由把 $V_\theta$ 在 $I_{w\lambda+\rho}$ 处局部化得到 [4, Ch.~IV, \S2, no.~5, Cor.~1 of Prop.~9]。这意味着
$$\Tor_p^{V_\theta}(\C_{\lambda+\rho},\C_{\lambda+\rho})=\Tor_p^{R_\lambda}(\C,\C),\qquad p\in\Z_+.$$
由于各 $R_{w\lambda}$ 彼此同构，
$$\operatorname{Card}W=\dim_\C V_\theta=\sum_{w\in W/W(\lambda)}\dim_\C R_{w\lambda}=\operatorname{Card}(W/W(\lambda))\dim_\C R_\lambda,$$
即 $\dim_\C R_\lambda=\operatorname{Card}W(\lambda)\ne 1$。因此，$R_\lambda$ 不是正则局部环，其同调维数无穷 [7, 17.3.1]，且对 $p\in\Z_+$ 有 $\Tor_p^{R_\lambda}(\C,\C)\ne 0$ [7, 17.2.11]。
\end{proof}

由此立即得到下述结果。

\begin{proposition}\label{prop:03.2.8}
设 $\theta$ 是 $\h^*$ 中由奇异元素构成的 Weyl 群轨道。则 $U_\theta$ 的同调维数无穷。
\end{proposition}

由 1.22 可推出下述结论。照旧置 $W(p)=\{w\in W\mid \ell(w)=p\}$。

\begin{proposition}\label{prop:03.2.9}
设 $F$ 是以 $\lambda$ 为最低权的有限维不可约 $\g$-模。则对任意 $p\in\Z_+$，
$$H_p(\n_x,F)=\sum_{w\in W(p)}\C_{w(\lambda-\rho)+\rho}.$$
\end{proposition}

\begin{proof}
显然 $\lambda-\rho$ 正则，故可应用 \ref{cor:03.2.6}。由 1.22 知，$F$ 的各阶局部化都是局部自由 $\calO_X$-模。因此，它们的高次几何纤维全为零，谱序列退化。公式立即由 1.22 得出。
\end{proof}

\begin{annote}
底本在 2.9 之前的过渡句与 2.9 的证明里都称「由 1.22」，但 1.22 只断言有限维 $\g$-模与有限生成 $\g$-模的张量仍有限生成，据此得不到「$F$ 的各阶局部化都是局部自由 $\calO_X$-模」，也得不到 2.9 的公式。给出这两件事的是 1.24（有限维不可约模的平移公式，取 $\mu=w(\lambda-\rho)$）。译文照录底本编号，读者请按 1.24 读。\hfill [来源：底本 p.46，对照 1.22 与 1.24 的陈述]
\end{annote}
